\documentclass[a4paper,12pt]{article}
\usepackage[utf8]{inputenc}
\usepackage[T1]{fontenc}
\usepackage[ngerman]{babel}
\usepackage{amsmath}
\usepackage{amssymb}
\usepackage{enumitem}
\usepackage{geometry}
\geometry{a4paper, margin=2.5cm}
\usepackage{parskip}

\title{Einsendeaufgaben -- Aufgabe 6.4\\[0.5em]
\large Modul 61112: Lineare Algebra}
\author{}
\date{}

\begin{document}
\maketitle

\section*{Aufgabe 6.4}

\begin{enumerate}[label=\alph*)]
  \item
  \begin{align*}
    \chi_A &= \det(X I_3 - A)
      = \det\begin{pmatrix} X-5 & -6 & -2 \\ 2 & X+2 & 1 \\ 0 & 0 & X-2 \end{pmatrix} \\
    &= (X-2)\big((X-5)(X+2)+12\big) \\
    &= (X-2)(X^2-3X+2) \\
    &= (X-2)^2(X-1)
  \end{align*}
  \[
    \Rightarrow \lambda_1 = 2 \text{ und } \lambda_2 = 1
  \]

  \item
  \begin{align*}
    \mathrm{Eig}_A(2) = \ker(A - 2I_3)
    &= \ker\begin{pmatrix} 3 & 6 & 2 \\ -2 & -4 & -1 \\ 0 & 0 & 0 \end{pmatrix} \\
    &= \ker\begin{pmatrix} 3 & 6 & 2 \\ 0 & 0 & 1 \\ 0 & 0 & 0 \end{pmatrix} \\
    &= \ker\begin{pmatrix} 1 & 2 & \tfrac{2}{3} \\ 0 & 0 & 0 \\ 0 & 0 & 1 \end{pmatrix} \\
    &= \left\langle \begin{pmatrix} 2 \\ -1 \\ 0 \end{pmatrix} \right\rangle
  \end{align*}

  \begin{align*}
    \mathrm{Eig}_A(1) = \ker(A - I_3)
    &= \ker\begin{pmatrix} 4 & 6 & 2 \\ -2 & -3 & -1 \\ 0 & 0 & 1 \end{pmatrix} \\
    &= \ker\begin{pmatrix} 2 & 3 & 1 \\ 0 & 0 & 0 \\ 0 & 0 & 1 \end{pmatrix} \\
    &= \ker\begin{pmatrix} 1 & \tfrac{3}{2} & 0 \\ 0 & 0 & 0 \\ 0 & 0 & 1 \end{pmatrix} \\
    &= \left\langle \begin{pmatrix} \tfrac{3}{2} \\ -1 \\ 0 \end{pmatrix} \right\rangle
     = \left\langle \begin{pmatrix} 3 \\ -2 \\ 0 \end{pmatrix} \right\rangle
  \end{align*}

  \item
  \[
    H_A(1) = \left\langle \begin{pmatrix} 3 \\ -2 \\ 0 \end{pmatrix} \right\rangle
  \]
  \begin{align*}
    H_A(2) = \ker\!\left( (A - 2I_3)^2 \right)
    &= \ker\begin{pmatrix} 3 & 6 & 2 \\ -2 & -4 & -1 \\ 0 & 0 & 0 \end{pmatrix}^2 \\
    &= \ker\begin{pmatrix} -3 & -6 & 0 \\ 2 & 4 & 0 \\ 0 & 0 & 0 \end{pmatrix} \\
    &= \ker\begin{pmatrix} 1 & 2 & 0 \\ 0 & 0 & 0 \\ 0 & 0 & 0 \end{pmatrix} \\
    &= \left\langle \begin{pmatrix} 2 \\ -1 \\ 0 \end{pmatrix}, \begin{pmatrix} 0 \\ 0 \\ 1 \end{pmatrix} \right\rangle
  \end{align*}

  $p = (1,1)$ ist die Rangpartition von $\mathrm{Eig}_A(2) \subsetneq \mathrm{Eig}_A^2(2)$.
  Die duale Partition von $p$ ist $q = (2)$.

  Wir wählen $v_{1,2} = \begin{pmatrix} 0 \\ 0 \\ 1 \end{pmatrix}
  \in \mathrm{Eig}_A^2(2) \setminus \mathrm{Eig}_A(2)$.

  Wir erhalten $v_{1,1} = (A - 2I_3)\, v_{1,2} = \begin{pmatrix} 2 \\ -1 \\ 0 \end{pmatrix}$.

  Wir können $S$ mit $B_1 = \left( \begin{pmatrix} 2 \\ -1 \\ 0 \end{pmatrix},
  \begin{pmatrix} 0 \\ 0 \\ 1 \end{pmatrix} \right)$ von $H_A(2)$ und
  $B_2 = \begin{pmatrix} 3 \\ -2 \\ 0 \end{pmatrix}$ von $H_A(1)$:
  \[
    S = \begin{pmatrix} 2 & 0 & 3 \\ -1 & 0 & -2 \\ 0 & 1 & 0 \end{pmatrix}
  \]

  Die Jordan'sche Normalform von $A$ ist
  \[
    \mathrm{diag}\big( J_2(2), J_1(1) \big)
    = \begin{pmatrix} 2 & 1 & 0 \\ 0 & 2 & 0 \\ 0 & 0 & 1 \end{pmatrix}
  \]
\end{enumerate}

\end{document}
